\begin{ouverture}[geometrie]{12}{Repérage et grandeurs en fonction d’une autre}{Partie 2 · Algèbre et grandeurs}
{Repère orthogonal · Coordonnées · Tableaux · Graphiques}
\objectif{Lire et placer des points dans un repère orthogonal.}
\objectif{Produire un tableau à partir d’une formule.}
\objectif{Lire une dépendance et reconnaître la proportionnalité.}
\prerequis{Relatifs et formules.}
\end{ouverture}
\pagestyle{otMain}\essentiel\fontsize{11}{15}\selectfont
\phantomsection\label{cours-12-1}
\essSection{1}{Deux coordonnées dans un ordre précis}
\begin{definition}Un repère orthogonal comporte deux axes perpendiculaires : l’axe horizontal des abscisses, l’axe vertical des ordonnées. Un point $A(3;-2)$ se repère par 3 sur le premier axe puis $-2$ sur le second. Les unités des axes peuvent être différentes.\end{definition}
\label{fig-cours-12-1-cours-repere}\figureGeo{cours-repere}

\begin{savoirfaire}[Lire et placer]Pour lire un point, projeter horizontalement et verticalement sur les axes. Pour le placer, partir de l’origine, suivre l’abscisse puis l’ordonnée. Sur l’axe horizontal, l’ordonnée est 0 ; sur l’axe vertical, l’abscisse est 0.\end{savoirfaire}
\begin{attention}$A(3;-2)$ et $B(-2;3)$ sont généralement différents. Les coordonnées sont un couple ordonné.\end{attention}
\clearpage
\phantomsection\label{cours-12-2}
\essSection{2}{Décrire une dépendance}
\begin{definition}Dire « le prix en fonction de la durée » signifie que la durée sert de donnée d’entrée et que le prix dépend de cette durée. On peut employer une formule, un tableau ou un graphique, sans étude générale des fonctions.\end{definition}
\begin{center}\renewcommand{\arraystretch}{1.35}\begin{tabular}{c|cccc}\toprule
Durée $t$ (h) & 0 & 1 & 2 & 4\\\midrule
Prix $3t+5$ (euros) & 5 & 8 & 11 & 17\\
\bottomrule\end{tabular}\end{center}
\label{fig-cours-12-2-cours-dependance}\figureGeo{cours-dependance}

\begin{savoirfaire}[Passer d’une représentation à l’autre]Lire les titres, les unités et les pas des graduations. Les points d’un tableau de données discrètes ne sont pas forcément à relier. Une formule permet de construire le tableau ; un graphique permet de visualiser les variations.\end{savoirfaire}
\begin{attention}Des points alignés ne suffisent pas pour une proportionnalité : la droite doit aussi passer par l’origine. Une courbe croissante n’est pas forcément une droite.\end{attention}

\clearpage\rubrique{Découverte}
\begin{blocLivre}{ex-12-1}\begin{exercice}[1]{Lire deux points}
Lire les coordonnées de A et B.
\label{fig-ex-12-1}\figureGeo{rep-1}
\end{exercice}\end{blocLivre}
\begin{blocLivre}{ex-12-2}\begin{exercice}[1]{Placer un point}
Placer $C(-2;-3)$ dans le repère.
\label{fig-ex-12-2}\figureGeo{rep-2}
\end{exercice}\end{blocLivre}
\begin{blocLivre}{ex-12-3}\begin{exercice}[1]{Sur les axes}
Placer $D(0;3)$ et $E(-4;0)$.
\label{fig-ex-12-3}\figureGeo{rep-3}
\end{exercice}\end{blocLivre}
\begin{blocLivre}{ex-12-4}\begin{exercice}[1]{Inverser les coordonnées}
Placer $F(1;3)$ et $G(3;1)$. Sont-ils confondus ?
\label{fig-ex-12-4}\figureGeo{rep-4}
\end{exercice}\end{blocLivre}
\begin{blocLivre}{ex-12-5}\begin{exercice}[1]{Lire une valeur de graphique}
La courbe donne une température en fonction de l’heure. Quelle température lit-on à 2 h ?
\label{fig-ex-12-5}\figureGeo{rep-5}
\end{exercice}\end{blocLivre}

\rubrique{Entraînement}
\begin{blocLivre}{ex-12-6}\begin{exercice}[2]{Construire un tableau}
Le prix en euros est $2t+4$, avec $t$ en heures. Donner les prix pour $t=0,1,3$ puis placer les points.
\label{fig-ex-12-6}\figureGeo{rep-6}
\end{exercice}\end{blocLivre}
\begin{blocLivre}{ex-12-7}\begin{exercice}[2]{Une courbe de carré}
Pour un carré de côté $c$ cm, l’aire est $c^2$. Compléter pour $c=0,1,2,3$, puis placer les points. Sont-ils alignés ?
\label{fig-ex-12-7}\figureGeo{rep-7}
\end{exercice}\end{blocLivre}
\begin{blocLivre}{ex-12-8}\begin{exercice}[2]{Coordonnées décimales}
Placer $H(-1{,}5;2{,}5)$ et $I(2{,}5;-0{,}5)$ sur le quadrillage gradué par demi-unités.
\label{fig-ex-12-8}\figureGeo{rep-8}
\end{exercice}\end{blocLivre}
\begin{blocLivre}{ex-12-9}\begin{exercice}[2]{Lire plusieurs moments}
Le graphique donne un volume d’eau. Lire le volume à 0, 2 et 5 minutes. Sur quelle période reste-t-il constant ?
\label{fig-ex-12-9}\figureGeo{rep-9}
\end{exercice}\end{blocLivre}
\begin{blocLivre}{ex-12-10}\begin{exercice}[2]{Des points alignés}
Les points du graphique sont alignés. La relation est-elle proportionnelle ? Justifier avec l’origine.
\label{fig-ex-12-10}\figureGeo{rep-10}
\end{exercice}\end{blocLivre}

\rubrique{Problèmes}
\begin{blocLivre}{ex-12-11}\begin{exercice}[3]{Un rectangle dans le plan}
A, B et C sont trois sommets consécutifs d’un rectangle : $A(-3;-1)$, $B(2;-1)$, $C(2;3)$. Placer le quatrième sommet D.
\label{fig-ex-12-11}\figureGeo{rep-11}
\end{exercice}\end{blocLivre}
\begin{blocLivre}{ex-12-12}\begin{exercice}[3]{Retrouver une formule}
Un parking coûte 5 euros fixes puis 2 euros par heure. Donner le prix en fonction de $t$, puis lire le prix pour 4 h sur le graphique.
\label{fig-ex-12-12}\figureGeo{rep-12}
\end{exercice}\end{blocLivre}
\begin{blocLivre}{ex-12-13}\begin{exercice}[3]{Comparer deux tarifs}
A coûte $3t$ euros ; B coûte $t+6$ euros. Pour les durées entières de 0 à 5 h, calculer les prix et repérer quand ils sont égaux.
\label{fig-ex-12-13}\figureGeo{rep-13}
\end{exercice}\end{blocLivre}
\begin{blocLivre}{ex-12-14}\begin{exercice}[3]{Lire une tendance sans extrapoler}
Le nuage de points donne des observations à 1, 2, 3 et 4 minutes. Peut-on affirmer la valeur à 5 minutes ? La relation observée est-elle proportionnelle ?
\label{fig-ex-12-14}\figureGeo{rep-14}
\end{exercice}\end{blocLivre}
\begin{blocLivre}{ex-12-15}\begin{exercice}[3]{Une réserve qui baisse}
Une réserve contient 24 L et perd régulièrement 3 L par minute, jusqu’à être vide. Produire un tableau pour 0, 2, 4, 6 et 8 minutes ; écrire une formule et interpréter sa limite d’utilisation.
\label{fig-ex-12-15}\figureGeo{rep-15}
\end{exercice}\end{blocLivre}

\clearpage
